\chapter*{Introducción}
Las aplicaciones distribuidas representan una respuesta integral a la necesidad emergente de sistemas capaces de gestionar procesos complejos, escalables y geográficamente dispersos en un mundo cada vez más globalizado y digitalizado. La dependencia de sistemas que operan efectiva y eficientemente a través de múltiples ubicaciones y zonas horarias se ha incrementado de manera significativa. En este contexto, se considera una realidad cotidiana la colaboración en tiempo real entre equipos de diferentes partes del mundo, un fenómeno facilitado en gran medida por las aplicaciones distribuidas.

Desde los inicios de las redes de computadoras, ha existido un deseo constante de compartir recursos, información y capacidades. A medida que la tecnología ha avanzado, las ambiciones en torno a la capacidad de procesar, analizar y actuar sobre los datos se han ampliado. Las aplicaciones distribuidas se erigen como el nexo entre estas ambiciones y la realidad tangible, permitiendo la colaboración entre múltiples sistemas y componentes, que operan como una entidad cohesiva, independientemente de las distancias físicas que los separan.

En este libro, se proporcionará una comprensión profunda de que las aplicaciones distribuidas no son simplemente una evolución de las aplicaciones monolíticas tradicionales. Por el contrario, representan una redefinición completa de cómo debe abordarse el diseño, desarrollo y mantenimiento del software en la actualidad. Se ha transitado de la centralización en un único servidor o sistema hacia un escenario en el que múltiples sistemas, a menudo distribuidos globalmente, cooperan armoniosamente.

El primer capítulo se dedicará a explorar la definición, características y relevancia actual de las aplicaciones distribuidas, proporcionando una visión integral que servirá como base para los capítulos subsiguientes. A través de ejemplos cotidianos, se ilustrará la omnipresencia de estas aplicaciones en diversos ámbitos de la vida diaria, desde las redes sociales hasta los sistemas bancarios que procesan transacciones en tiempo real.

Es importante señalar que, a pesar de las numerosas ventajas que ofrecen las aplicaciones distribuidas, también presentan desafíos únicos. La complejidad intrínseca de diseñar sistemas que operan en múltiples ubicaciones, junto con la necesidad de asegurar la confidencialidad y privacidad de los datos, así como gestionar la escalabilidad y el rendimiento, son algunos de los aspectos que se abordarán en detalle en este libro.

Por ende, este libro busca no solo proporcionar una comprensión teórica, sino también práctica, de cómo se conceptualizan, implementan y mantienen las aplicaciones distribuidas. A través de cada página, se aspira a acercar al lector un paso más hacia la conversión en un experto en el campo, capaz de comprender y contribuir al paisaje en constante evolución de la Ingeniería de Software.

Por consiguiente, se invita al lector a emprender este viaje hacia el fascinante universo de las aplicaciones distribuidas. Bienvenido a un mundo donde la colaboración, la innovación y la tecnología convergen para crear soluciones que transforman la realidad tal como la conocemos.